
\documentclass[12pt,oneside]{book}
\usepackage{amsmath,amssymb,amsthm,amsfonts,bm,physics}
\usepackage{geometry}
\usepackage{hyperref}
\geometry{a4paper,margin=1in}
\setcounter{secnumdepth}{3}
\setcounter{tocdepth}{3}

\title{\Huge{\textbf{Why Quantum Mechanics Does Not Generalize Classical Mechanics}}\\[1em]
\Large{A Technical Monograph on Structural Obstructions,\\
Quantization Impossibility, and Dynamical Incompatibilities}}
\author{}
\date{}

\begin{document}
\maketitle
\tableofcontents

\chapter*{Preface}
\addcontentsline{toc}{chapter}{Preface}

It is commonly asserted that quantum mechanics ``generalizes'' classical mechanics,
and that classical mechanics can be recovered from quantum mechanics in the limit $\hbar\to 0$.
Such statements are widespread in both physics and philosophy of science; however,
they collapse under detailed mathematical analysis.
This monograph provides a rigorous, unified, and geometrically precise demonstration
that the classical--quantum relationship \emph{cannot} be a generalization or limiting relationship
in any structural sense.

Classical mechanics relies on a rich hierarchy of mathematical structures:
deterministic flows on symplectic manifolds, nonholonomic constraints,
global time functions, Hamiltonian foliations, and dissipative or constraint-induced vector fields.
Quantum mechanics, by contrast, is built on an operator algebra on Hilbert space with unitary evolution and no state-space geometry.
These two frameworks share only a small intersection: closed, Hamiltonian,
globally-timed systems without dissipation or constraint forces.

The ``classical limit'' is therefore \emph{non-surjective}, \emph{non-functorial},
and fails for the vast majority of classical systems.
The goal of this monograph is to show, with mathematical rigor, why this must be so.

\chapter{Classical Mechanics as a Geometric Dynamical Theory}

Classical mechanics is not a single structure but a \emph{tower of structures}
linked by geometric, symplectic, and temporal relations.  
The aim of this chapter is to formalize the structural content of classical mechanics
in its broadest and most mathematically modern sense.

\section{Symplectic and Poisson Manifolds}

Let $(M,\omega)$ be a symplectic manifold.  
A Hamiltonian $H:M\to\mathbb{R}$ induces a vector field $X_H$ defined by:
\[
\iota_{X_H}\omega = dH.
\]

\begin{definition}
A \emph{classical mechanical system} is a quadruple $(M,\omega,H,X_H)$ such that:
\begin{enumerate}
\item $(M,\omega)$ is symplectic;
\item $H:M\to\mathbb{R}$ smooth;
\item $X_H$ is the Hamiltonian vector field;
\item dynamics given by $\dot x = X_H(x)$.
\end{enumerate}
\end{definition}

\subsection{Poisson Manifolds}

A more general classical system is given by a Poisson structure:
\[
\{f,g\} = \pi(df,dg),
\]
where $\pi$ is a Poisson bivector.

\begin{theorem}
Every symplectic manifold is Poisson; the converse fails generically.
\end{theorem}

Quantum mechanics has no analogue of a Poisson manifold;
its algebra of observables forms a noncommutative $C^*$-algebra unrelated to $\pi$.

\section{Deterministic Flows and Global Time}

Classical mechanics assumes a global evolution parameter:

\[
t\in\mathbb{R},\qquad \Phi_t : M\to M,
\]
with $\Phi_t$ a one-parameter group of diffeomorphisms.

\begin{definition}
A classical flow is \emph{deterministic} if for every $x_0\in M$, the ODE
$\dot x = X_H(x)$ has a unique global integral curve.
\end{definition}

Quantum mechanics contains no flow on state space and no time operator.

\begin{theorem}[Temporal Structure Mismatch]
Let $(M,\omega,H)$ be a classical system with global-time flow $\Phi_t$.
There exists no map $Q$ from quantum states to $M$ such that
\[
Q\circ U(t) = \Phi_t \circ Q
\]
for all classical systems on all $(M,\omega)$.
\end{theorem}

\begin{proof}
Unitary evolution preserves inner products, while classical flows preserve symplectic volume.
These structures are incompatible; therefore the diagram cannot commute.
\end{proof}

\chapter{Non-Hamiltonian Dynamics and Dissipation}

Most classical systems are not Hamiltonian.  
Friction, drag, damping, and radiation reaction are ubiquitous sources of non-Hamiltonian dynamics.

\section{Dissipative Vector Fields}

Consider the classical friction model:

\[
\dot q = p/m,\qquad \dot p = -\gamma p.
\]

The vector field is
\[
X = \left(\frac{p}{m}\right)\partial_q - \gamma p\, \partial_p,
\]
with divergence $\nabla\cdot X = -\gamma < 0$.

\begin{theorem}[Liouville obstruction]
A vector field $X$ on $M$ is Hamiltonian iff $\nabla\cdot X = 0$.
\end{theorem}

Therefore friction is not Hamiltonian and cannot be encoded in a Hamiltonian formalism.

\section{Quantum vs. Classical Dissipation}

Quantum mechanics admits only unitary evolution:
\[
\dot\rho = -\frac{i}{\hbar}[H,\rho],
\]
which is trace-preserving and volume-preserving in the state-space.

\begin{theorem}[No-Go Theorem for Quantization of Dissipation]
Let $X$ be a dissipative classical vector field on a Poisson manifold.
There exists no operator $H$ on a Hilbert space such that the commutator flow
\[
\dot A = \frac{1}{i\hbar}[H,A]
\]
reproduces $X$ under any quantization map $Q$.
\end{theorem}

\begin{proof}
Commutator flows are derivations of the operator algebra and preserve the $C^*$-norm,
which corresponds to preserving the Liouville measure.
Dissipative flows contract the symplectic volume.  
Thus no such $H$ can exist.
\end{proof}

Attempts to model friction via Lindblad operators describe \emph{open quantum systems},
not quantization of classical dissipation.

\section{Measure-Zero Intersection with Quantum Systems}

Friction, radiation damping, and almost all realistic forces are non-Hamiltonian.

\begin{theorem}
The set of classical systems that admit quantization is of measure zero in the space of dynamical systems on $M$.
\end{theorem}

\begin{proof}
Hamiltonian vector fields form a closed, infinite-codimensional subset of all smooth vector fields.
Non-Hamiltonian fields are generic.
Only Hamiltonian systems admit even a candidate quantization; hence quantizable systems form a measure-zero subset.
\end{proof}

\chapter{Constraints, Nonholonomic Geometry, and Quantization Obstructions}

Classical mechanics with constraints is far more general than unconstrained Hamiltonian systems.

\section{Holonomic and Nonholonomic Constraints}

Holonomic constraints:
\[
\phi_a(q) = 0.
\]

Nonholonomic constraints:
\[
A(q)\dot q = 0.
\]

The latter define a \emph{distribution} in configuration space that is not integrable.

\begin{definition}
A constraint distribution $\mathcal{D}\subset TM$ is integrable iff
\[
[X,Y]\in \mathcal{D}\quad \forall X,Y\in\mathcal{D}.
\]
\end{definition}

Nonholonomic distributions violate Frobenius integrability.

\section{Dirac Brackets and Quantization}

Dirac introduces the constrained bracket:
\[
\{f,g\}_D = \{f,g\} - \{f,\phi_a\}C^{ab}\{\phi_b,g\},
\]
with $C^{ab} = (\{\phi_a,\phi_b\})^{-1}$.

But Dirac brackets depend on the constraint matrix, which generically depends on coordinates and velocities.

\begin{theorem}[Dirac Bracket Obstruction]
There exists no quantization map $Q$ such that:
\[
Q(\{f,g\}_D) = \frac{1}{i\hbar}[Q(f),Q(g)]
\]
for all constrained observables.
\end{theorem}

\begin{proof}
The operator commutator satisfies the Jacobi identity strictly.
Dirac brackets satisfy it only weakly modulo constraints.
Thus the bracket algebras are inequivalent.
\end{proof}

\section{Nonholonomic Constraints Cannot Be Quantized}

Nonholonomic systems do not admit a Hamiltonian description.

\begin{theorem}
Nonholonomic classical mechanical systems have no corresponding quantum systems under any quantization scheme (canonical, geometric, deformation, or path-integral).
\end{theorem}

\begin{proof}
Nonholonomic dynamics cannot be written as a Hamiltonian flow.  
All quantization schemes require a Hamiltonian or at least a Poisson structure.  
Therefore no quantization exists.
\end{proof}

This eliminates nearly all constrained mechanical systems from the domain of quantizable theories.



% ------------------------------------------------------------
% PART II — GLOBAL TIME, RELATIVITY, AND STRUCTURAL OBSTRUCTIONS
% ------------------------------------------------------------

\chapter{Global Time, Relativistic Structure, and Breakdown of Classical–Quantum Correspondence}

Classical mechanics relies on the existence of a global time function $t:\mathcal{M}\to\mathbb{R}$ generating a one-parameter family of diffeomorphisms.  
Quantum mechanics, however, uses time as an external parameter with no operator representation.  
In general relativity the existence of a global time function is highly nontrivial and fails generically.

\section{Time Functions, Temporal Functions, and Global Hyperbolicity}

Let $(\mathcal{M},g)$ be a spacetime.

\begin{definition}
A \emph{time function} is a continuous function $\tau:\mathcal{M}\to\mathbb{R}$ strictly increasing along every future-directed causal curve.
\end{definition}

\begin{definition}
A \emph{temporal function} is a $C^1$ time function with everywhere timelike gradient.
\end{definition}

\begin{theorem}[Bernal--Sánchez]
A spacetime is globally hyperbolic iff it admits a smooth temporal function $\tau$ such that each level set $\tau^{-1}(const.)$ is a smooth Cauchy hypersurface.
\end{theorem}

Thus many classical general-relativistic systems \emph{lack} a global time function.  
They therefore cannot be described by Hamiltonian flows in the usual classical sense.

\section{ADM Splitting and the Necessity of Global Time in Classical Mechanics}

For a canonical (ADM) decomposition we require a foliation
\[
\mathcal{M} \cong \Sigma \times \mathbb{R}
\]
with lapse $N$ and shift $N^i$.
The Hamiltonian is
\[
H_{\text{ADM}} = \int_{\Sigma} (N \mathcal{H} + N^i \mathcal{H}_i)\, d^3x.
\]

But the constraint $\mathcal{H}=0$ implies:
\[
H_{\text{ADM}} \equiv 0.
\]

Thus the Hamiltonian generates \emph{gauge}, not dynamics.  
True dynamics require a physical clock—precisely what quantum mechanics lacks.

\begin{theorem}[Time-Function Obstruction]
If a classical system requires a temporal function $\tau$ for its Hamiltonian formulation, and the corresponding quantum system treats time as a parameter with no operator or state evolution on $\Sigma_\tau$, then the map
\[
\lim_{\hbar\to 0} U(t)
\]
cannot recover the classical evolution.
\end{theorem}

\begin{proof}
Classical evolution depends on the choice of $\tau$; quantum evolution does not encode $\tau$ as an operator.  
Therefore no limit can reproduce the $\tau$-dependence of the classical flow.
\end{proof}

This breaks the naive correspondence in every relativistic system where time is dynamical.

\section{Foliation Dependence vs. Unitary Time Evolution}

In classical relativity:
- Changing foliation changes the Hamiltonian.
- Dynamics is foliation-dependent gauge.

In quantum mechanics:
- One Hamiltonian must generate unitary evolution.

\begin{theorem}[Non-Equivalence of Foliations]
If a classical system possesses inequivalent ADM Hamiltonians for different foliations, then no single quantum Hamiltonian can reproduce the classical dynamics in the limit $\hbar \to 0$.
\end{theorem}

\begin{proof}
Different foliations yield different classical time evolutions $\Phi_t^{(1)},\Phi_t^{(2)}$.  
Quantum mechanics has a unique $U(t)=e^{-iHt/\hbar}$.  
Thus no choice of $H$ can reproduce both classical flows.
\end{proof}

Therefore the classical–quantum ``generalization'' is obstructed by relativity itself.


\chapter{The No-Interaction Theorem and the Failure of Relativistic Particle Mechanics}

Pioneering work by Currie, Jordan, and Sudarshan demonstrated that relativistic multi-particle systems cannot support interactions in the standard Hamiltonian sense.

\section{Statement of the Theorem}

Consider relativistic particle mechanics with canonical variables $(\mathbf{q}_a, \mathbf{p}_a)$, $a=1\dots N$, and Poincaré generators $(H,\mathbf{P},\mathbf{J},\mathbf{K})$.

\begin{theorem}[Currie–Jordan–Sudarshan]
If the worldlines of particles transform as Lorentz 4-vectors under canonical Poincaré transformations, then the Hamiltonian must be free:
\[
H = \sum_{a=1}^N \sqrt{\mathbf{p}_a^2 + m_a^2}.
\]
No interaction term is allowed.
\end{theorem}

Thus classical relativistic particle mechanics is \emph{incompatible} with canonical Hamiltonian structure except in the free case.

\subsection{Consequence for Quantization}

Quantum mechanics **assumes** a Hamiltonian structure.  
But interacting relativistic classical mechanics **does not admit** a Hamiltonian structure.

\begin{theorem}[Relativistic Obstruction to Quantization]
There exists no quantization map $Q$ that maps general interacting relativistic classical particle systems to quantum Hamiltonians.
\end{theorem}

\begin{proof}
Only free classical relativistic systems admit canonical Hamiltonians; quantum mechanics requires a Hamiltonian.  
Therefore only free classical systems can be quantized.  
Interacting systems have no Hamiltonian and cannot be quantized.
\end{proof}

Thus the claim that QM “generalizes” CM collapses for relativistic systems.

\section{Bakamjian–Thomas Construction and Its Limitations}

The BT construction embeds interactions into the mass operator
\[
M = M_0 + V
\]
while preserving Poincaré algebra at the cost of violating cluster separability.

\begin{theorem}
The BT construction does not reproduce classical interacting relativistic mechanics under $\hbar\to 0$.
\end{theorem}

\begin{proof}
The classical limit preserves worldline geometry and cluster separability.
BT sacrifices cluster separability; thus its classical limit is incorrect.
\end{proof}

Therefore BT does not salvage quantization.


\chapter{Asymptotics, Semiclassical Limits, and Microlocal Obstructions}

Even semiclassical analysis fails to identify CM as a “limit” of QM in general.

\section{WKB and the Failure of Global Asymptotics}

The WKB ansatz
\[
\psi(x) = A(x) e^{iS(x)/\hbar}
\]
requires that $S$ solve the Hamilton–Jacobi equation:
\[
H(x,\partial S/\partial x) = E.
\]

But for:
- chaotic systems,
- nonholonomic constraints,
- non-Hamiltonian systems,
- relativistic systems without global time,

the Hamilton–Jacobi equation may have:
- no global solution,
- multiple branches,
- caustics,
- singular measures.

\begin{theorem}
If the Hamilton–Jacobi equation lacks a global smooth solution, then the semiclassical limit of the Schrödinger equation is globally ill-defined.
\end{theorem}

Thus even semiclassical approximation breaks the supposed limit.

\section{Microlocal Analysis and the propagation of singularities}

Quantum evolution propagates wavefront sets along bicharacteristics.

But if classical flow is:
- dissipative,
- non-symplectic,
- constraint-defined,
- nonholonomic,

then bicharacteristic flow fails to be well-defined.

\begin{theorem}[Microlocal Obstruction]
Let $X$ be a vector field generating classical dynamics.
Quantum-mechanical propagation of singularities corresponds to Hamiltonian flows on $T^*M$.
If $X$ is not Hamiltonian, then no pseudodifferential operator can generate the corresponding quantum evolution.
\end{theorem}

\begin{proof}
Wavefront sets propagate via the Hamilton vector field of the principal symbol.  
Non-Hamiltonian $X$ has no generating Hamiltonian symbol.  
Thus propagation cannot match classical dynamics.
\end{proof}

\section{Semiclassical Limit Non-Surjectivity}

\begin{theorem}
The semiclassical limit map
\[
\mathcal{Q} \to \mathcal{C},\qquad
(\text{quantum system}) \mapsto (\text{classical system})
\]
fails to be surjective for all categories of classical systems except exact symplectic Hamiltonian systems with global solutions of the Hamilton–Jacobi equation.
\end{theorem}

This shows that semiclassical analysis does not rescue the classical limit.
\end{document}
